% LSPV-4 REPORT::MODELBRIDGE-0026::2026-10-12T14:00:46Z
% !TEX program = xelatex
\documentclass[11pt]{article}
\usepackage[margin=0.86in]{geometry}
\usepackage{fontspec,amsmath,amssymb,amsthm,booktabs,enumitem,xurl,hyperref}
\setlength{\emergencystretch}{3em}
\newtheorem{theorem}{Theorem}
\newtheorem{corollary}[theorem]{Corollary}
\newcommand{\EXACT}{\textbf{PROVED/EXACT}}
\newcommand{\KNOWN}{\textbf{KNOWN THEOREM}}
\newcommand{\COMP}{\textbf{COMPUTATIONAL}}
\newcommand{\OPEN}{\textbf{OPEN}}
\newcommand{\RETIRED}{\textbf{FALSIFIED/RETIRED}}
\title{MODELBRIDGE-0026:\\Fixed Normalized-Time Cut Exposure Collapses under Polynomial Pre-Delay}
\author{LSPV-4 --- Universal Core Bridge}
\date{12 October 2026}
\begin{document}
\maketitle

\begin{abstract}
This run falsifies one precise attempted repair of the arbitrary-algorithm
clean-core bridge.  RT-0009 already showed that every fixed polynomial trace
prefix can be hidden behind a higher-degree delay.  Here the checkpoint is
allowed to scale with the wrapped machine's exact total runtime: for any fixed
rational $0<\alpha<1$ and any total uniform deterministic polynomial-time
decider $A$, there is an extensionally equivalent polynomial-time decider
$W_{A,\alpha}$ whose mutable configuration at time
$\lfloor\alpha Q(N)\rfloor$, where $Q(N)$ is its exact total runtime on every
$N$-bit input, depends only on $N$.  The read-only input is deliberately not
erased; treating its mere presence as ``exposure'' is an external source route,
not solver-paid state.  Thus a nonconstant clean-core compatibility object
$H(F)$ cannot be forced into mutable state at any one fixed normalized-time
cut.  A $1{,}533{,}168$-check finite audit passed.  Exogenous online challenges,
peak-space obligations, all-cuts statements, and a model-matched endpoint
remain open.  P versus NP remains open.
\end{abstract}

\section{Executive status and exact target}

\begin{center}\small
\begin{tabular}{p{0.38\linewidth}p{0.20\linewidth}p{0.34\linewidth}}
\toprule Claim & Status & Verdict\\ \midrule
\texttt{NORMALIZED-CUT-PREDELAY1} & \EXACT &
Every fixed rational fraction of exact total runtime can be placed inside an
input-oblivious polynomial prelude.\\
\texttt{FIXED-FRACTION-CUT-EXPOSURE1} & \RETIRED &
No nonconstant source object is forced into mutable solver state at one such
cut for every extensionally equivalent polynomial decider.\\
\texttt{EXOGENOUS-CHALLENGE-COMPILER1} & \OPEN &
No compiler from every static SAT decider to an irrevocable externally timed
challenge interface is known.\\
\texttt{PEAK-OR-ALL-CUTS-TRANSFER1} & \OPEN &
The theorem says nothing about maximum workspace or every time in an interval.\\
P versus NP & \OPEN & No unrestricted deterministic SAT lower bound follows.\\
\bottomrule
\end{tabular}
\end{center}

Let the clean core be the canonical duplicate-free proper signed
Monotone-$3$SAT-$(2,2)$ family used by current root \texttt{main.tex}.  For a
formula $F$, let $H(F)$ denote any proposed nonconstant semantic compatibility
object (residual service, boundary table, accepted-node label, or other exact
source functional).  Input length $N$ is the complete encoded bit length.

The attacked target has the quantifier form
\[
 \exists\alpha\in(0,1)\cap\mathbb Q\ \forall A\in\mathsf{Dec}_{P}
 \ \forall^{\infty}F:\quad
 H(F)\preceq \pi\!\operatorname{Conf}_A
 \bigl(F,\lfloor\alpha Q_A(|F|)\rfloor\bigr),             \tag{1}
\]
where $A$ has an input-length-exact runtime $Q_A$, $\pi$ projects to mutable
control/work/head state, and $\preceq$ means exact extraction with all costs
charged.  A statement using the read-only input itself instead of $\pi$ is
source extraction, not evidence that the execution acquired or retained
$H(F)$.  The SAME\_REALIZATION ledger includes the wrapper, length scan, exact
clock, child simulation, retained answer bit, dummy steps, mutable tapes, head
positions, extraction, and any source accesses.

\section{Normalized-cut pre-delay theorem}

\begin{theorem}[\texttt{NORMALIZED-CUT-PREDELAY1}]
Fix a language $L\subseteq\{0,1\}^*$, a total uniform deterministic decider
$A$ for $L$, an integer-valued time-constructible polynomial bound $S(N)$ for
an exact-clock simulation of $A$, and a rational
$\alpha=a/b\in(0,1)$ in lowest terms.  There is a total uniform deterministic
decider $W_{A,\alpha}$ and an exact polynomial runtime $Q$ such that
\begin{enumerate}[label=(\roman*)]
\item $W_{A,\alpha}(x)=A(x)$ for every $x$;
\item $W_{A,\alpha}$ runs for exactly $Q(N)$ steps on every $N$-bit input;
\item for all $x,x'\in\{0,1\}^N$,
\[
 \pi\!\operatorname{Conf}_{W_{A,\alpha}}
    (x,\lfloor\alpha Q(N)\rfloor)
 =
 \pi\!\operatorname{Conf}_{W_{A,\alpha}}
    (x',\lfloor\alpha Q(N)\rfloor),                     \tag{2}
\]
where $\pi$ excludes read-only input contents; and
\item if $S(N)=O(N^K)$ then $Q(N)=O(N^K)$, with constants allowed to depend on
$A$ and $\alpha$.
\end{enumerate}
\end{theorem}

\begin{proof}
The wrapper first performs a symbol-oblivious length scan: its transition on
input symbols $0$ and $1$ is identical, and only the endmarker changes phase.
It resets the input head and enters an input-oblivious dummy prelude.  Absorb
the scan into that prelude.  Let
\[
 D(N)=L(N)+\left\lceil\frac{aS(N)+b}{b-a}\right\rceil,
 \qquad Q(N)=D(N)+S(N).                                  \tag{3}
\]
where $L(N)$ is the exact symbol-oblivious length-scan/reset time.  After
exactly $D(N)$ prelude steps the wrapper simulates $A$ under the exact
$S(N)$ clock.  When $A$ halts, the wrapper retains its one-bit answer and idles
until the clock expires, then returns that bit.  This proves totality,
uniformity, semantic equivalence, and exact runtime.

Equation (3) gives
\[
 (b-a)D(N)\ge aS(N)+b,
\]
and hence $a(D(N)+S(N))\le bD(N)-b$.  Therefore
\[
 \lfloor\alpha Q(N)\rfloor <D(N).                       \tag{4}
\]
At that time the child has not begun and the prelude's mutable configuration
is a deterministic function of $N$, $A$, and $\alpha$, but not of the input
contents.  Equality (2) follows.  Finally $D(N)=O(S(N))$, so the wrapper
preserves polynomiality and the exponent $K$ up to constants.
\end{proof}

\begin{corollary}[Failure of a fixed normalized cut]
If $H$ is nonconstant on infinitely many equal-length input pairs, no theorem
of form (1) can hold for all extensionally equivalent polynomial deciders.
\end{corollary}
\begin{proof}
Choose equal-length $x,x'$ with $H(x)\ne H(x')$ and apply Theorem~1 to any
polynomial decider.  Its wrapped mutable cut states are identical, so no exact
function of that state distinguishes $H(x)$ from $H(x')$.
\end{proof}

This result is compatible with $P=NP$.  Under that hypothesis, apply the
construction to a polynomial-time clean-core or general SAT decider; the
wrapped solver is still polynomial.  There is no contradiction with any
unconditional lower bound.  In particular, the wrapper may increase constants
and exact runtime while preserving the unknown exponent; FLMV-style
time--space bounds apply only when one same-model realization satisfies their
simultaneous endpoint.

\section{What the theorem does not say}

The theorem does not show that no input-dependent state ever occurs.  It does
not lower-bound or upper-bound peak workspace.  It does not defeat a cut chosen
adaptively after observing the computation, an assertion quantified over all
times in a long interval, or an externally delivered challenge that arrives
after an independently verified irrevocable commitment.  It also does not
remove the source from the read-only input tape.  These distinctions are
load-bearing:
\begin{itemize}[leftmargin=*]
\item including the source in the ``state'' makes exposure trivial but supplies
no solver-paid acquisition theorem;
\item waiting until the child begins imports the arbitrary implementation
schedule;
\item demanding a commitment/challenge interface changes the computational
model and requires a total compiler from every static decider; and
\item demanding every cut or peak state needs a separate machine lower bound,
not semantic-state counting.
\end{itemize}

Thus the smallest repair is \texttt{EXOGENOUS-CHALLENGE-COMPILER1}: specify a
causal interface whose challenge time is not a function of the wrapped
program's self-declared runtime, prove that every uniform polynomial SAT
decider can be compiled to it with all preprocessing, rereading, recomputation,
storage, and exponent costs charged, and then match an unconditional lower
bound to that same compiled realization.  No such theorem is verified.

\section{Reproducible bounded audit}

The dependency-free invocation
\begin{verbatim}
timeout 30s python experiments/review/normalized_cut_predelay_audit.py \
 --seed 202610121400 \
 --output results/raw/review/2026-10-12_normalized_cut_predelay.json
\end{verbatim}
completed in $1.467387$ seconds; Python bytecode compilation passed.

\begin{center}\small
\begin{tabular}{lr}\toprule Check & Count\\ \midrule
Exact fractional-cut inequalities & 6,048\\
Semantic answer preservation & 762,048\\
Hostile family/encoding/wrapper cells & 762,048\\
Same-length mutable-cut collisions & 3,024\\ \midrule
Total & 1,533,168\\ \bottomrule
\end{tabular}
\end{center}

The grid uses lengths $2$--$64$, child exponents $1$--$8$, fractions
$1/10,1/4,1/2,3/4,9/10,99/100$, seven families (clean core, Horn, 2SAT,
affine/Tseitin, bounded treewidth, independent products, parity), six encoding
variants, and DIRECT/DEAD/COPY/DELAY/INTERP/RECOMP wrappers.  Script SHA-256:
\nolinkurl{dbb535e3e8670a33a77e59a8f39e0f17783a93ac46baffb64ef238fe852a9680}.
Raw JSON SHA-256:
\nolinkurl{059b81cfbae65cbe152ae285bfa25cfa29d78811a5a08fc2b32ef9e51cedf39d}.
The grid is \COMP; Theorem~1, not finite extrapolation, proves the symbolic
statement.  No invocation exceeded the repository's 30-second rule, so no
chunking was required.

\section{Primary-literature and transfer audit}

Hartmanis--Stearns supplies the foundational time-bounded machine framework,
and Blum emphasizes that complexity is attached to program/machine
realizations rather than extension alone \cite{hs,blum}.  The present wrapper
theorem is elementary and proved above; it does not invoke Blum speedup.
Hennie--Stearns gives a constructive model simulation with explicit overhead,
which reinforces the requirement to charge a compiler rather than identify
models by analogy \cite{hennie}.  FLMV gives genuine SAT time--space lower
bounds on its stated machine model, but not a compiler from arbitrary
polynomial SAT code into a causal cut protocol \cite{flmv}.

\begin{description}[leftmargin=0cm,labelwidth=0cm]
\item[] \textbf{DIRECTLY TRANSFERABLE THEOREM.}
The exact-clock and dummy-prelude construction transfers to deterministic
multitape/RAM-like models with time-constructible polynomial clocks; every
transition and work cell is charged.  The same construction works relative to
an oracle.
\item[] \textbf{USEFUL ANALOGY REQUIRING A NEW TRANSFER THEOREM.}
Online chronograms, communication cuts, control-theoretic observability
windows, VLSI clocking, and interactive proof challenges impose exogenous
interfaces.  They become relevant only after a compiler preserves their
causal hypotheses for every static SAT decider.
\item[] \textbf{MISLEADING ANALOGY / FAILED TRANSFER.}
A physical clock edge, a chosen trace fraction, a communication round, or an
observable latch does not constrain an unrestricted static SAT program by
itself.
\end{description}

The theorem relativizes verbatim, so it cannot supply the nonrelativizing
hinge blocked by Baker--Gill--Solovay \cite{bgs}.  Giving the wrapper algebraic
oracle access does not invalidate it; hence it also does not cross
algebrization \cite{aw}.  It proposes no large constructive circuit property,
so Natural Proofs is not the immediate obstruction.  Schur squares, quadratic
relation ranks, Hilbert functions, syzygies, and collision entropy are
source-side invariants and are \textbf{FAILED TRANSFER} for forcing mutable
state at this cut.  LWE-style black-box extraction and special-soundness
collision extraction are \textbf{ANALOGY/CANDIDATE}; the missing unconditional
static-to-causal compiler remains.  Cryptographic binding and random-oracle
assumptions are inadmissible.

\section{Corpus reconciliation}

\subsection*{CORPUS REVIEWED}
The current GitHub tree was surveyed at 2,589 paths.  Root \texttt{main.tex},
\texttt{AGENTS.md}, \texttt{CODEX.md}, the complete PNP research skill, actual
\texttt{main4.tex}/\texttt{main5.tex}, all current task specifications, and
verified MAIN \texttt{SYNTHESIS-0023} were read.  The run read the newest
relevant specialist reports HALTCANON1, PRODSTREAM1, WRAPORBIT1,
MODELBRIDGE-0025, and TIGHTDELAY1; RT-0009 and its fixed-checkpoint script/raw
result; relevant historical MAIN and model-bridge reports; governance; and the
HEC/VSC-HEC/G-PPC-HEC/CRS-PPC/CCSPC/RCCN and random-challenge surfaces.
Repository search covered normalized cuts, fractional checkpoints, delay,
wrapper covariance, exact clocks, exposure, and source-defined cuts.  Primary
papers on program-indexed complexity, model simulation, SAT time--space bounds,
relativization, and algebrization were checked before promotion.

The cyclic predecessor is LSPV-3.  Its newest verified canonical artifact is
\path{research/agents/lspv-3/2026-10-12_1237_LSPV-3_WRAPORBIT1.tex}; no artifact
for its next scheduled cycle was visible.  The exact missing scientific body
and tool error were unavailable, so the failure is classified
\textbf{UNKNOWN / incomplete-write evidence}.  A nonblocking same-body-only
recovery request was written and verified at
\path{research/review/outbox/model_bridge/2026-10-12_1400_LSPV3_WRITE_RECOVERY.md}.
No mathematical content was reconstructed.

\subsection*{DUPLICATION / CONTRADICTION AUDIT}
RT-0009 already proves that a fixed $N^r$ trace prefix is defeated by an
$N^{r+1}$ delay.  HALTCANON1 already retires final retained-state necessity,
and TIGHTDELAY1 already charges delay exponent.  This report does not rename
those results.  The new increment is the exact inequality (3)--(4) for a cut
defined as a fixed fraction of the wrapper's own exact runtime.  Repository
search found no earlier normalized-time theorem.  The result is consistent
with WRAPORBIT1's failure of finite wrapper certification and PRODSTREAM1's
sequential-reuse control.

\subsection*{WHAT CHANGED RELATIVE TO ROOT MAIN}
Root MAIN already marks the arbitrary-algorithm transfer and causal compiler
\OPEN.  The LSPV-4 region is sharpened to record that neither a fixed absolute
prefix nor one fixed normalized-time cut can force source-dependent mutable
state.  The highest-value remaining bridge is consequently an exogenous
challenge/commitment compiler or a peak/all-cuts theorem with a same-model
endpoint.  No executive section or other specialist region is changed.

Actual \texttt{main4.tex} remains a conditional ARC framework, and actual
\texttt{main5.tex} remains a conditional wrapper/charged-contraction
framework.  Neither contains a proved all-decider compiler.  They are not
silently redefined by this report.

\subsection*{ADOPTED / REJECTED / RETIRED PRIOR COMPONENTS}
Adopted are RT-0009's fixed-prefix theorem, TIGHTDELAY1's charged delay,
HALTCANON1's halt cleanup, MODELBRIDGE-0025's external adapter requirement,
SAME\_REALIZATION, exact encoding length, and all hostile controls.  Rejected
is a fixed normalized cut as a wrapper-invariant exposure point.  Retired is
\texttt{FIXED-FRACTION-CUT-EXPOSURE1}.  Exogenous challenge cuts, peak space,
all-cuts obligations, and source-only external extraction are not retired.

\subsection*{STALE OR OVERSTATED CLAIMS FOUND}
Any claim that ``fraction of total runtime'' removes the DELAY wrapper is
overstated: the wrapper can make its prelude occupy more than any fixed
fraction while staying polynomial.  Any claim counting the unchanged
read-only source as solver-retained compatibility state conflates external
source access with paid internal acquisition.

\section{Single next target and handoff}

The unique bridge target is \texttt{EXOGENOUS-CHALLENGE-COMPILER1}.  State an
interface in which (i) the source and program encodings are explicit, (ii) a
commitment is certified before an externally chosen challenge, (iii) the
challenge cannot be rescheduled by DEAD/COPY/DELAY/INTERP/RECOMP, (iv) all
preprocessing, rereading, reconstruction, storage, and arbitrary polynomial
exponents are charged, (v) the same compiled realization enters an
unconditional lower bound, and (vi) the proof identifies a nonrelativizing and
nonalgebrizing step.  A source-only construction must instead prove a
$K$-independent charged contraction and cannot charge its object to the
original solver's state.

\begin{thebibliography}{9}
\bibitem{hs} J.~Hartmanis and R.~E.~Stearns,
\emph{On the Computational Complexity of Algorithms}, Trans. AMS 117 (1965),
285--306. \url{https://doi.org/10.1090/S0002-9947-1965-0170805-7}.
\bibitem{blum} M.~Blum,
\emph{A Machine-Independent Theory of the Complexity of Recursive Functions},
JACM 14(2) (1967), 322--336.
\url{https://doi.org/10.1145/321386.321395}.
\bibitem{hennie} F.~C.~Hennie and R.~E.~Stearns,
\emph{Two-Tape Simulation of Multitape Turing Machines}, JACM 13(4) (1966),
533--546. \url{https://doi.org/10.1145/321356.321362}.
\bibitem{flmv} L.~Fortnow, R.~Lipton, D.~van Melkebeek, A.~Viglas,
\emph{Time--Space Lower Bounds for Satisfiability}, JACM 52(6) (2005),
835--865. \url{https://doi.org/10.1145/1101821.1101822}.
\bibitem{bgs} T.~Baker, J.~Gill, R.~Solovay,
\emph{Relativizations of the P=?NP Question}, SIAM J. Comput. 4(4) (1975),
431--442. \url{https://doi.org/10.1137/0204037}.
\bibitem{aw} S.~Aaronson and A.~Wigderson,
\emph{Algebrization: A New Barrier in Complexity Theory}, ACM TOCT 1(1)
(2009), Article 2. \url{https://eccc.weizmann.ac.il/report/2008/005/}.
\end{thebibliography}

\paragraph{Final status and verification sentinel.}
\texttt{NORMALIZED-CUT-PREDELAY1}: \EXACT.  Finite audit: \COMP.\\
\texttt{FIXED-FRACTION-CUT-EXPOSURE1}: \RETIRED.\\
\texttt{EXOGENOUS-CHALLENGE-COMPILER1},
\texttt{PEAK-OR-ALL-CUTS-TRANSFER1}, a model-matched unconditional endpoint,
nonrelativization, nonalgebrization, and P versus NP: \OPEN.
P versus NP remains open.
% END_OF_VERIFIED_REPORT::LSPV-4::MODELBRIDGE-0026
\end{document}
